\section{Gallery, with settings}
\label{app:gallery}

Every figure in this paper is reproducible from a short parameter list, because a
layer is a parameter list --- there is no asset, no baked texture, and no captured
raster anywhere in the pipeline. This appendix gives the settings for the teaser and
for two further compositions, in the same units the interface uses: world units for
lengths, degrees for layer rotation, radians for the per-member rotation
$\theta$.

\subsection{Teaser (Figure~\ref{fig:teaser})}

Five panels, each $620\times956$ device pixels at zoom $1.5$ on white ---
so $413\times637$ world units --- stroke thickness $1.5$ world units throughout,
$3\times3$ supersampling. At the printed panel width of $1.36$\,in that is $456$\,dpi
and a stroke a third of a point wide, which is the finest hairline that survives
the press; it is also what lets each panel carry sixty members across instead of
thirty. Both families of a pair are the same ink, \texttt{\#14171b}.

\begin{center}\footnotesize
\begin{tabular}{@{}llll@{}}
  \toprule
  Panel & Family A & Family B & Fringes \\
  \midrule
  1 combs   & lines, $s{=}6.5$, $\alpha{=}0.44$ & lines, $s{=}6.5$, $\alpha{=}0.50$ & lines \\
  2 rings   & circles, $s{=}8$ at $x{=}{-}90$ & circles, $s{=}8$ at $x{=}90$ & hyperbolae \\
  3 walking & triangles, $s{=}3.5$, $\varphi{=}2$, & --- & envelopes \\
            & $\delta{=}(1.2,0)$, $\theta{=}0.03$ & & \\
  4 spirals & spiral, $s{=}5$, $B{=}90$ & the same, mirrored in $y$ & rays \\
  5 flow    & lines, $s{=}5$ & the same $+$ \textsc{swirl}, $a{=}3.5$ & streamlines \\
  \bottomrule
\end{tabular}
\end{center}

\noindent Here $\alpha$ is the direction of $\nabla\ph$, in radians.

Four observations, since this is the first thing a reader sees.

Panel~3 is a \emph{single} family, not a pair. Its fringes come from its own
members running near-tangent to one another, which happens because $\theta\ne0$
turns each successive triangle a little further than the last. The classical
two-grating account has nothing to say about this configuration; the
index-difference account handles it without modification, because $D$ is a
difference of index fields and nothing requires the two fields to come from
different layers. The dark curves sweeping through that panel are the family's
\emph{envelope} --- the locus where consecutive members touch. They are the same
object as Lemma~\ref{lem:envelope}, seen rather than bounded, and they are the
reason the sufficient interval of Theorem~\ref{thm:window} has to be as wide as it
is: near an envelope the nearest index changes fast. Its $\rad(-\delta)/s$ is
close to but below $1$, which by \S\ref{sec:limits} puts it in the wide-interval
band near the degenerate locus. This is not an accident of taste. The settings that
look most alive are the ones nearest the case where the guarantee runs out, which
is the practical reason the bound had to be exact rather than merely safe.

Panel~4 is worth a line of arithmetic. Reflecting an Archimedean spiral negates
the $\theta$ term, so $D=-2\hat b\arg(p)/s$ with $\hat b=Ms/2\pi$ the used
pitch --- a function of the
polar angle alone: the fringes are exactly $2M$ rays, here thirty-six. Note what
this predicts and the picture confirms --- \emph{rotating} a spiral produces no
fringes at all, since a rotation only adds a constant to $D$. That is the
zoetrope property, and it falls out of the theorem rather than needing a separate
argument.

\subsection{Catalog plate (Figure~\ref{fig:plate})}

Settings for all twenty-six panels are in the plate script
(Appendix~\ref{app:repro}). Two conventions are worth stating: every panel covers
$500$ world units square at stroke thickness $1.5$ in those units --- twice the
extent a screen-sized view would show, because a beat needs several of its periods
in frame before it reads as a beat --- and no panel is a crop, scale or resample of
another. The perturbed copy in the lower half of a cell is a second field
evaluation, not a transformed image of the upper half.

\subsection{Contour compositions (Figures~\ref{fig:contourproof}--\ref{fig:instrument})}

Every panel is two \textsc{parallel} layers of pitch $s=\contourCarrier$, alike in
everything but the field carried by the second. So the settings are the expression,
the gain $a$, and the extent that maps world units into its $x$ and $y$ --- three
values, listed below with what the measurement of \S\ref{sec:beyond} found for
each. Every expression here is one of the presets of Table~\ref{tab:exprs}, so a
row can be reproduced by choosing it from the editor and typing two numbers. \emph{Finest} is the narrowest
on-screen contour interval $c/\lvert\nabla f\rvert$ in pixels, \emph{resolved} the
fraction of probes above the limit \eqref{eq:contournyquist}, and the errors are
displacements of the fringe peak from the true contour in pixels, against an
estimator floor of $\contourFloor$ / $\contourFloorP$\,px.

% GENERATED by paper/tools/numbers.mjs
\begin{center}\footnotesize
\begin{tabular}{@{}lrrrrrr@{}}
  \toprule
  Field & $a$ & Extent & $c=1/a$ & Finest & Resolved & Error mean / p95 \\
  \midrule
  \textsc{saddle} & 3 & 200 & 0.333 & 15.6 & 98\% & 0.29 / 0.56 \\
  \textsc{dipole} & 8 & 210 & 0.125 & 8.1 & 82\% & 0.22 / 0.55 \\
  \textsc{bumps} & 3 & 220 & 0.333 & 26.8 & 100\% & 0.80 / 4.21 \\
  \textsc{swirl} & 1.2 & 200 & 0.833 & 48.4 & 100\% & 0.51 / 1.59 \\
  \textsc{ripple} & 2.4 & 400 & 0.417 & 26.5 & 100\% & 0.26 / 0.78 \\
  \textsc{terrain} & 2.2 & 320 & 0.455 & 56.5 & 100\% & 0.53 / 1.49 \\
  \bottomrule
\end{tabular}
\end{center}


The dipole row is the one we would point at. Its field is by far the steepest, so
it is the panel where \eqref{eq:contournyquist} binds hardest --- its contours
close to $\contourDipoleFinest$\,px apart, and only $\contourDipoleResolved\%$ of
the frame resolves --- and it is also the panel that most obviously looks like a
physics figure. Both facts follow from the same gradient.

The remaining panels differ only in framing. Figure~\ref{fig:contourproof} uses
\textsc{saddle} at $a=0.1$ and extent $50$, a much coarser interval so the bands
are wide enough to read on paper and the true curves have room to be drawn over
them. Figure~\ref{fig:instrument} uses \textsc{swirl} at $a=2$, extent $180$, and
\textsc{bumps} at $a=2.6$, extent $200$, each shown twice --- once as rendered and
once under the envelope view, whose only parameter is a contrast of $2.4$ about the
stack's nominal coverage.
